\section{Protobuf Compatibility Definitions}

\begin{definition}[Encoding Field]
	An encoding field is a tuple $(id, tag, value)$ which encodes the field
	number, tag and value of one field of a proto descriptor.
\end{definition}

Just like how the message descriptor is a sequence of field descriptors, the
encoding type of a message is an ordered list of encoding fields.

\begin{definition}[Field Compatibility]
	For fields $f_1 = (id_1, tag_1, v_1)$ and $f_2 = (id_2, tag_2, v_2)$, $f_2$
	with descriptor $d_2$ is a compatible update to $f_1$ with descriptor $d_1$ if
	$id_1 = id_2$ and for all values $v_1$, $serialize_{d_1}\ f_1 = bs$ and there
	exists a $v_2$ such that $parse_{d_2}\ bs = v_2$ and $v_1 \prec v_2$ according to
	some relation $\prec$.
\end{definition}

An example value relation is given in Section~\ref{sec:val-rel}.

\begin{definition}[Message Compatibility]
	A protobuf message descriptor $d_2$ is a compatible update to $d_1$ if every
	field in $d_1$ has a compatible field in $d_2$.
\end{definition}

Formally speaking, the message compatibility relation ($\preceq$) is given in
Section~\ref{sec:comp-rel}.

\begin{definition}[Tag Function]
  The tag function $tag : \text{Type} \rightarrow \text{Tag}$ maps the \texttt{proto}
  level type to the corresponding tag as listed in Table~\ref{tab:tags}.
\end{definition}

It is also important to understand the default values of each protobuf
type. When a field is serialized with the default value or not set in the
message, the parser places the default value in the resulting struct. The
primary difference between an implicit field and an optional field is that the
optional field can differentiate between the serialized field being omitted from
the message or being present, but having the default value.

\begin{definition}[Default Function]
  The default function
  $default: (t:\text{proto type}) \rightarrow \llbracket t \rrbracket$ returns the default
  value for each type.
  
  \[ default(t) = \left\{ 
    \begin{array}{l@{\quad:\quad}l}
      0 & t \in \left\{\begin{array}{c}
        \mathtt{int32}, \mathtt{int64}, \mathtt{uint32},
        \mathtt{uint64}, \mathtt{sint32}, \mathtt{sint64}, \\ \mathtt{fixed32},
        \mathtt{fixed64}, \mathtt{sfixed32}, \mathtt{sfixed64}, \mathtt{enum}
      \end{array}\right\} \\
      \mathtt{false} & t \text{ is } \mathtt{bool} \\
      ``" & t \text{ is } \mathtt{string} \\
      % 0 width space needed so that the brackets don't parse as a distance to
      % space out the elements of the array... for some unknown reason.
      \hspace{0pt}[\;] & t \text{ is } \mathtt{bytes}
    \end{array}\right.
  \]
\end{definition}

% Local Variables:
% citar-bibliography: ("../pollux.bib")
% TeX-master: "../pollux.tex"
% jinx-local-words: "protobuf"
% End:
